\section{A fixed-shift law on the exact coprime shell}
\label{sec:fixed-shell}

We now prove \eqref{eq:intro-fixed-shell-law}.  The arithmetic function
\begin{equation}
 d^*(r)=2^{\omega(r)}
 =\#\{(a,b)\in\N^2:ab=r,\ (a,b)=1\}
 \label{eq:unitary-divisor-function}
\end{equation}
counts unitary divisors.  We use $a\parallel r$ to mean
$a\mid r$ and $(a,r/a)=1$.

\begin{lemma}[Exact shell coefficient]
\label{lem:exact-shell-coefficient}
For every $r\geq1$,
\begin{equation}
 \sum_{\substack{a\parallel r\\a>1}}u(a)
 =\log\rad(r)-d^*(r)+1.
 \label{eq:exact-shell-coefficient}
\end{equation}
Consequently
\begin{equation}
 \cK_h(X)=
 \sum_{r\geq1}
 \bigl(\log\rad(r)-d^*(r)+1\bigr)
 \eta(r+h)W(r/X).
 \label{eq:shell-product-ledger}
\end{equation}
\end{lemma}

\begin{proof}
The unitary divisors of
$r=\prod_pp^{\nu_p}$ are obtained by choosing, for each $p\mid r$, either
$1$ or the complete prime-power block $p^{\nu_p}$.  Hence there are
$2^{\omega(r)}$ of them.  The von Mangoldt weight is nonzero on a unitary
divisor precisely when that divisor is one prime-power block, and therefore
\[
 \sum_{a\parallel r}\Lambda(a)
 =\sum_{p\mid r}\log p=\log\rad(r).
\]
Since $u(1)=-1$, deleting $a=1$ adds $1$ and proves
\eqref{eq:exact-shell-coefficient}.  Grouping \eqref{eq:intro-k-shell} by
$r=ab$ proves \eqref{eq:shell-product-ledger}.
\end{proof}

The leading term comes entirely from the mismatch between the ordinary
mean of $d^*$ and its mean at a shifted prime.

\begin{theorem}[Unitary-divisor Titchmarsh input]
\label{thm:unitary-titchmarsh-input}
For every fixed nonzero integer $h$,
\begin{equation}
 \sum_{p\leq y}d^*(p-h)
 =\beta_hy+O_h\left(\frac y{\log y}\right),
 \label{eq:unitary-titchmarsh-input}
\end{equation}
where $d^*(n)=0$ for $n\leq0$ and $\beta_h$ is given by
\eqref{eq:beta-h}.
\end{theorem}

This is the $k=2$ case of Corollary~1.2 of
\citet{BiaoWang2026}.  In the notation of that paper, the displayed
constant initially contains
\[
 \frac{\zeta(2)\zeta(3)}{\zeta(6)}
 \prod_p\left(1-\frac1{p^2-p+1}\right).
\]
The local identity
\[
 \frac{1-p^{-6}}{(1-p^{-2})(1-p^{-3})}
 \left(1-\frac1{p^2-p+1}\right)=1
\]
reduces it to \eqref{eq:beta-h}.

\begin{lemma}[The two unitary-divisor means]
\label{lem:two-unitary-means}
For fixed $h\ne0$,
\begin{align}
 \sum_{r\geq1}d^*(r)W(r/X)
 &=\frac{XI_0(W)}{\zeta(2)}\log X+O_W(X),
 \label{eq:ordinary-unitary-mean}\\
 \sum_{r\geq1}d^*(r)\Lambda(r+h)W(r/X)
 &=\beta_hXI_0(W)\log X+O_{h,W}(X).
 \label{eq:prime-unitary-mean}
\end{align}
\end{lemma}

\begin{proof}
The Dirichlet series identity
\begin{equation}
 \sum_{r\geq1}\frac{d^*(r)}{r^s}
 =\frac{\zeta(s)^2}{\zeta(2s)}
 \qquad(\Re s>1)
 \label{eq:unitary-dirichlet-series}
\end{equation}
has a double pole at $s=1$ with leading coefficient $1/\zeta(2)$.
Standard Selberg--Delange summation, or the convolution identity obtained
from \eqref{eq:unitary-dirichlet-series}, gives
\[
 \sum_{r\leq y}d^*(r)
 =\frac y{\zeta(2)}\log y+O(y).
\]
Partial summation against $W(r/X)$ proves
\eqref{eq:ordinary-unitary-mean}.  Only its $X\log X$ coefficient is used;
all secondary $X$ terms are absorbed in the remainder.

For \eqref{eq:prime-unitary-mean}, first restrict $r+h$ to primes and put
$p=r+h$.  Partial summation in \eqref{eq:unitary-titchmarsh-input}, with
the test function
\[
 y\longmapsto \log y\,W((y-h)/X),
\]
gives
\[
 \sum_p d^*(p-h)\log p\,W((p-h)/X)
 =\beta_hXI_0(W)\log X+O_{h,W}(X).
\]
Here the main integral is
\[
 \int_0^\infty \log y\,W((y-h)/X)\,dy
 =XI_0(W)\log X+O_{h,W}(X).
\]
The targets $r+h=p^j$, $j\geq2$, contribute $O_{h,W}(X^{3/4})$, say:
there are $O(X^{1/2})$ such targets, and
$d^*(p^j-h)\ll_\delta X^\delta$.  Taking $\delta=1/8$ and absorbing the
remaining logarithm proves the displayed bound and hence the full von
Mangoldt form.
\end{proof}

It remains to show that the logarithmic-radical part has no $X\log X$
term.

\begin{lemma}[Radical residual]
\label{lem:radical-residual}
For every fixed nonzero $h$,
\begin{equation}
 \sum_{r\geq1}\log\rad(r)\eta(r+h)W(r/X)
 =O_{h,W}(X).
 \label{eq:radical-residual}
\end{equation}
Also
\begin{equation}
 \sum_{r\geq1}\eta(r+h)W(r/X)=O_{h,W}(X/\log X).
 \label{eq:constant-residual-bound}
\end{equation}
\end{lemma}

\begin{proof}
Write
\begin{equation}
 L(r)=\log r-\log\rad(r)
 =\sum_{\substack{p^j\mid r\\j\geq2}}\log p.
 \label{eq:squarefull-log-correction}
\end{equation}
The classical zero-free-region prime number theorem, after translating the
variable by fixed $h$ and applying smooth partial summation, gives
\begin{equation}
 \sum_r\log r\,\eta(r+h)W(r/X)=O_{h,W}(X).
 \label{eq:log-residual-bound}
\end{equation}
The weaker error shown here is more than sufficient.

We claim
\begin{equation}
 \sum_{r\asymp_WX}L(r)(1+\Lambda(r+h))\ll_{h,W}X.
 \label{eq:squarefull-correction-bound}
\end{equation}
For the unweighted part, insert \eqref{eq:squarefull-log-correction}; the
convergent sum
$\sum_{p,j\geq2}(\log p)/p^j$ gives $O(X)$, and the endpoint cost is
$O(X^{1/2}\log X)$.

For the von Mangoldt part, first remove proper prime-power targets.  Since
$L(r)\ll_W\log X$ on the support,
\[
 \sum_{\substack{\ell^\nu\asymp_W X\\\nu\geq2}}
 (\log\ell)L(\ell^\nu-h)
 \ll_{h,W}X^{1/2}\log X.
\]
It remains to consider prime targets.  First take $p^j\leq X^{1/2}$.  If
$p\nmid h$, the Brun--Titchmarsh inequality in the reduced class
$h\pmod{p^j}$ gives \citep{IwaniecKowalski2004}
\[
 \sum_{\substack{r\asymp_WX\\p^j\mid r}}
 \one_{r+h\ {\rm prime}}\Lambda(r+h)
 \ll_{h,W}\frac X{\varphi(p^j)}.
\]
Indeed,
\[
 \sum_p\sum_{j\geq2}\frac{\log p}{\varphi(p^j)}
 =\sum_p\frac{\log p}{(p-1)^2}<\infty.
\]
If $p\mid h$, a prime target divisible by $p$ would have to equal $p$ and
is absent from the support for large $X$.  For $p^j>X^{1/2}$, use
$\Lambda\leq\log(3X)$ and count the $O(X/p^j+1)$ multiples.  Explicitly,
\[
 \log X
 \sum_{\substack{j\geq2,\ p^j>X^{1/2}\\p^j\ll_W X}}
 (\log p)\left(\frac X{p^j}+1\right)
 \ll_W X^{3/4}\log X;
\]
the squares dominate, while $j\geq3$ is smaller.  This proves
\eqref{eq:squarefull-correction-bound}.  Combining it with
\eqref{eq:log-residual-bound} proves \eqref{eq:radical-residual}.

Finally, \eqref{eq:constant-residual-bound} is the ordinary smooth prime
number theorem, again after translating by fixed $h$.
\end{proof}

\begin{theorem}[Fixed-shift coprime-shell law]
\label{thm:fixed-shell-law}
Let $h\ne0$ be fixed.  Then
\begin{equation}
 \cK_h(X)=
 \left(\frac1{\zeta(2)}-\beta_h\right)
 XI_0(W)\log X+O_{h,W}(X).
 \label{eq:fixed-shell-law}
\end{equation}
For $h=2$, the leading coefficient is
$6/\pi^2-1/3$.
\end{theorem}

\begin{proof}
Insert \eqref{eq:shell-product-ledger}.  The contribution of
$-d^*(r)\eta(r+h)$ is the difference of
\eqref{eq:ordinary-unitary-mean} and
\eqref{eq:prime-unitary-mean}.  The radical and constant terms are
$O(X)$ by \cref{lem:radical-residual}.
\end{proof}

\begin{remark}[Where the fixed shift was bought]
The theorem does not estimate each orientation $b$ separately.  The exact
sum over all $b$ produces $d^*(r)$, and the unitary-divisor Titchmarsh
theorem evaluates that aggregate at shifted primes.  If one retains only
$b=1$, the sum becomes
\[
 \sum_a(\Lambda(a)-1)(\Lambda(a+h)-1)W(a/X),
\]
which is the fixed-shift prime-pair covariance.  Thus
\cref{thm:fixed-shell-law} is genuine fixed-$h$ progress at the radial-shell
level, not an endpoint estimate in disguise.
\end{remark}
